% ===========================================================
%  RESEARCH ABSTRACT
%  Self-contained fragment — use % ===========================================================
%  RESEARCH ABSTRACT
%  Self-contained fragment — use % ===========================================================
%  RESEARCH ABSTRACT
%  Self-contained fragment — use % ===========================================================
%  RESEARCH ABSTRACT
%  Self-contained fragment — use \input{Abstract} in main doc,
%  or uncomment the preamble below for standalone compilation.
% ===========================================================

% --- Standalone preamble (uncomment to compile independently) ---
\documentclass[12pt, a4paper]{article}
\usepackage[T1]{fontenc}
\usepackage[utf8]{inputenc}
\usepackage[english]{babel}
\usepackage{geometry}
\geometry{margin=2.5cm}
\usepackage{setspace}
\usepackage{microtype}
\usepackage{siunitx}
\begin{document}
% ----------------------------------------------------------------

\begin{center}
  {\large\bfseries
    Learning Motor Control Strategies in Musculoskeletal Systems\\[4pt]
    Using Deep Reinforcement Learning:\\[2pt]
    Replacing Physics-Based Muscle Solvers with Trained Policies}\\[10pt]
  {\normalsize Mohamed Naceur Hammami 
  %\quad|\quad
   %Supervisor: Prof.\ Taysir Rezgui
   }\\[4pt]
  %{\small\itshape
   %Master's in Complex and Intelligent Systems ---
   %\'{E}cole Polytechnique de Tunisie, University of Carthage, 2026}
\end{center}

\bigskip
\noindent\rule{\textwidth}{0.6pt}
\smallskip
\begin{center}{\large\bfseries Abstract}\end{center}
\smallskip
\noindent\rule{\textwidth}{0.6pt}
\bigskip

% ---- PROBLEM STATEMENT ----------------------------------------
\noindent
Real-time musculoskeletal control requires solving the Hill fiber--tendon
equilibrium at each simulation timestep to obtain per-muscle forces. The
standard Newton--Raphson approach iterates 5--40 times per muscle, imposing
runtime computation that scales with both the number of muscles and the
control frequency --- a meaningful burden for resource-constrained embedded
rehabilitation systems. This work proposes replacing the online iterative
Newton--Raphson solver with a trained deep reinforcement learning (DRL)
policy that shifts computation entirely to an offline training phase,
enabling control via a single neural network forward pass at deployment;
the policy outputs steady-state muscle activations $a \in [0,1]^9$, which
are then passed directly to the MuJoCo physics engine.

% ---- METHOD ---------------------------------------------------
\medskip\noindent
A nine-muscle upper-limb model was implemented in MuJoCo~3.2 with a
Hill-type muscle representation including tendon compliance, pennation
dynamics, and first-order activation dynamics, with anatomical parameters
drawn from Holzbaur et~al.\ (2005) and Winter (2009). Four DRL algorithms
--- DDPG, TD3, SAC, and PPO --- were trained over $2 \times 10^{6}$ steps
with domain randomization ($\pm15$\,\% segment mass, $\pm20$\,\% actuator
gain, $\pm30$\,\% joint damping) across ten independent random seeds with
Optuna-optimized hyperparameters. Statistical significance and effect sizes
were assessed using non-parametric testing and bootstrap confidence intervals.

% ---- COMPUTATIONAL PERFORMANCE --------------------------------
\medskip\noindent
\textbf{Computational performance.}\;
Under identical framework conditions (Python/PyTorch on a host CPU), the
trained SAC policy executes in $295\,\mu\text{s}$ per control step versus
$2165\,\mu\text{s}$ for the Newton--Raphson equilibrium solver. This figure is
reported for completeness but does \emph{not} support a deployment claim: it
times forward simulation rather than load sharing, and against a direct solve
of the load-sharing problem the network is roughly an order of magnitude
slower. Critically, the policy's execution time
is \emph{deterministic and branch-free} --- it consists of two fixed matrix
multiplications regardless of the biomechanical state --- whereas the N-R
solver's loop count varies with convergence conditions, precluding hard
real-time guarantees on embedded hardware. This property makes the policy
directly deployable on neural processing units (NPUs). The parameter count
($393{,}750$) and memory footprint (\SI{1.5}{\mega\byte} single precision,
\SI{385}{\kilo\byte} at eight-bit quantisation) are likewise fixed and
independent of the biomechanical state.

% ---- TASK PERFORMANCE -----------------------------------------
\medskip\noindent
\textbf{Task performance.}\;
On a three-dimensional reaching task, the best configuration attained a mean
final error of \SI{0.106}{\meter} and a closest approach of
\SI{0.074}{\meter}, ending within \SI{5}{\centi\meter} in $34$\,\% of episodes
and within \SI{10}{\centi\meter} in $58$\,\%. Under a strict success criterion
--- within \SI{2}{\centi\meter} and stationary --- performance was
$0$--$4$\,\%, a criterion that a privileged controller with full state access
satisfies only $\approx 7$\,\% of the time on this model. No statistical
significance testing is reported: the comparison uses three seeds and the
principal analysis a single seed.

% ---- MUSCLE FORCE FIDELITY ------------------------------------
\medskip\noindent
\textbf{Load-sharing fidelity.}\;
The learned per-muscle forces were compared against classical static
optimisation, posed on the identical trajectory with the identical required
joint torques and verified to machine precision. The comparison is conditional:
static optimisation is given the policy's own movement and asked only how to
distribute force across it, so trajectory realism is not assessed. The policy reproduces the
coordination pattern of the classical solution ($r = 0.81 \pm 0.04$ over five
independent training runs; pattern cosine $0.813$ against a shuffled null of
$0.372$) but expends $1.91 \pm 0.15$ times the minimum effort needed to generate
those torques --- about 44\,\% above the $1.33$ a conventional controller
attains on the same task. That discrepancy proved to be an artefact of the
reward's weighting: the effort term carried roughly one twenty-eighth the weight
of the precision bonus, and raising it brings the ratio to $1.34$, at the
achievable floor, with agreement rising to $r = 0.932$. The excess is anatomically localised: the
three deltoids agree to within $2.5$--$10$\,\% normalised RMSE, whereas the
lateral and medial heads of triceps are driven at three to four and a half
times the prescribed force, indicating sustained co-contraction against the
biceps group. Across all fifteen trained policies the discrepancy ranges from
$1.21$ to $2.30\times$, and the ordering by muscular economy is approximately
the inverse of the ordering by reaching accuracy.

% ---- CONCLUSION -----------------------------------------------
\medskip\noindent
These results support replacing the Newton--Raphson solver with a trained
policy where the quantity of interest is which muscles act and in what
relative proportion, and do not support it where the absolute magnitude of
individual muscle forces matters. Amortising solver computation into a
one-time training phase is computationally viable; the learned solution
recovers the structure of the classical load-sharing criterion without ever
being shown it, but not its efficiency.

\bigskip
\noindent\textbf{Keywords:}
musculoskeletal model,
Hill muscle model,
Newton--Raphson solver,
deep reinforcement learning,
Soft Actor-Critic,
motor control,
muscle synergies,
domain randomization,
amortized computation,
embedded rehabilitation,
MuJoCo.

% --- End standalone document (uncomment if compiling independently) ---
\end{document}
% ----------------------------------------------------------------------
 in main doc,
%  or uncomment the preamble below for standalone compilation.
% ===========================================================

% --- Standalone preamble (uncomment to compile independently) ---
\documentclass[12pt, a4paper]{article}
\usepackage[T1]{fontenc}
\usepackage[utf8]{inputenc}
\usepackage[english]{babel}
\usepackage{geometry}
\geometry{margin=2.5cm}
\usepackage{setspace}
\usepackage{microtype}
\usepackage{siunitx}
\begin{document}
% ----------------------------------------------------------------

\begin{center}
  {\large\bfseries
    Learning Motor Control Strategies in Musculoskeletal Systems\\[4pt]
    Using Deep Reinforcement Learning:\\[2pt]
    Replacing Physics-Based Muscle Solvers with Trained Policies}\\[10pt]
  {\normalsize Mohamed Naceur Hammami 
  %\quad|\quad
   %Supervisor: Prof.\ Taysir Rezgui
   }\\[4pt]
  %{\small\itshape
   %Master's in Complex and Intelligent Systems ---
   %\'{E}cole Polytechnique de Tunisie, University of Carthage, 2026}
\end{center}

\bigskip
\noindent\rule{\textwidth}{0.6pt}
\smallskip
\begin{center}{\large\bfseries Abstract}\end{center}
\smallskip
\noindent\rule{\textwidth}{0.6pt}
\bigskip

% ---- PROBLEM STATEMENT ----------------------------------------
\noindent
Real-time musculoskeletal control requires solving the Hill fiber--tendon
equilibrium at each simulation timestep to obtain per-muscle forces. The
standard Newton--Raphson approach iterates 5--40 times per muscle, imposing
runtime computation that scales with both the number of muscles and the
control frequency --- a meaningful burden for resource-constrained embedded
rehabilitation systems. This work proposes replacing the online iterative
Newton--Raphson solver with a trained deep reinforcement learning (DRL)
policy that shifts computation entirely to an offline training phase,
enabling control via a single neural network forward pass at deployment;
the policy outputs steady-state muscle activations $a \in [0,1]^9$, which
are then passed directly to the MuJoCo physics engine.

% ---- METHOD ---------------------------------------------------
\medskip\noindent
A nine-muscle upper-limb model was implemented in MuJoCo~3.2 with a
Hill-type muscle representation including tendon compliance, pennation
dynamics, and first-order activation dynamics, with anatomical parameters
drawn from Holzbaur et~al.\ (2005) and Winter (2009). Four DRL algorithms
--- DDPG, TD3, SAC, and PPO --- were trained over $2 \times 10^{6}$ steps
with domain randomization ($\pm15$\,\% segment mass, $\pm20$\,\% actuator
gain, $\pm30$\,\% joint damping) across ten independent random seeds with
Optuna-optimized hyperparameters. Statistical significance and effect sizes
were assessed using non-parametric testing and bootstrap confidence intervals.

% ---- COMPUTATIONAL PERFORMANCE --------------------------------
\medskip\noindent
\textbf{Computational performance.}\;
Under identical framework conditions (Python/PyTorch on a host CPU), the
trained SAC policy executes in $295\,\mu\text{s}$ per control step versus
$2165\,\mu\text{s}$ for the Newton--Raphson equilibrium solver. This figure is
reported for completeness but does \emph{not} support a deployment claim: it
times forward simulation rather than load sharing, and against a direct solve
of the load-sharing problem the network is roughly an order of magnitude
slower. Critically, the policy's execution time
is \emph{deterministic and branch-free} --- it consists of two fixed matrix
multiplications regardless of the biomechanical state --- whereas the N-R
solver's loop count varies with convergence conditions, precluding hard
real-time guarantees on embedded hardware. This property makes the policy
directly deployable on neural processing units (NPUs). The parameter count
($393{,}750$) and memory footprint (\SI{1.5}{\mega\byte} single precision,
\SI{385}{\kilo\byte} at eight-bit quantisation) are likewise fixed and
independent of the biomechanical state.

% ---- TASK PERFORMANCE -----------------------------------------
\medskip\noindent
\textbf{Task performance.}\;
On a three-dimensional reaching task, the best configuration attained a mean
final error of \SI{0.106}{\meter} and a closest approach of
\SI{0.074}{\meter}, ending within \SI{5}{\centi\meter} in $34$\,\% of episodes
and within \SI{10}{\centi\meter} in $58$\,\%. Under a strict success criterion
--- within \SI{2}{\centi\meter} and stationary --- performance was
$0$--$4$\,\%, a criterion that a privileged controller with full state access
satisfies only $\approx 7$\,\% of the time on this model. No statistical
significance testing is reported: the comparison uses three seeds and the
principal analysis a single seed.

% ---- MUSCLE FORCE FIDELITY ------------------------------------
\medskip\noindent
\textbf{Load-sharing fidelity.}\;
The learned per-muscle forces were compared against classical static
optimisation, posed on the identical trajectory with the identical required
joint torques and verified to machine precision. The comparison is conditional:
static optimisation is given the policy's own movement and asked only how to
distribute force across it, so trajectory realism is not assessed. The policy reproduces the
coordination pattern of the classical solution ($r = 0.81 \pm 0.04$ over five
independent training runs; pattern cosine $0.813$ against a shuffled null of
$0.372$) but expends $1.91 \pm 0.15$ times the minimum effort needed to generate
those torques --- about 44\,\% above the $1.33$ a conventional controller
attains on the same task. That discrepancy proved to be an artefact of the
reward's weighting: the effort term carried roughly one twenty-eighth the weight
of the precision bonus, and raising it brings the ratio to $1.34$, at the
achievable floor, with agreement rising to $r = 0.932$. The excess is anatomically localised: the
three deltoids agree to within $2.5$--$10$\,\% normalised RMSE, whereas the
lateral and medial heads of triceps are driven at three to four and a half
times the prescribed force, indicating sustained co-contraction against the
biceps group. Across all fifteen trained policies the discrepancy ranges from
$1.21$ to $2.30\times$, and the ordering by muscular economy is approximately
the inverse of the ordering by reaching accuracy.

% ---- CONCLUSION -----------------------------------------------
\medskip\noindent
These results support replacing the Newton--Raphson solver with a trained
policy where the quantity of interest is which muscles act and in what
relative proportion, and do not support it where the absolute magnitude of
individual muscle forces matters. Amortising solver computation into a
one-time training phase is computationally viable; the learned solution
recovers the structure of the classical load-sharing criterion without ever
being shown it, but not its efficiency.

\bigskip
\noindent\textbf{Keywords:}
musculoskeletal model,
Hill muscle model,
Newton--Raphson solver,
deep reinforcement learning,
Soft Actor-Critic,
motor control,
muscle synergies,
domain randomization,
amortized computation,
embedded rehabilitation,
MuJoCo.

% --- End standalone document (uncomment if compiling independently) ---
\end{document}
% ----------------------------------------------------------------------
 in main doc,
%  or uncomment the preamble below for standalone compilation.
% ===========================================================

% --- Standalone preamble (uncomment to compile independently) ---
\documentclass[12pt, a4paper]{article}
\usepackage[T1]{fontenc}
\usepackage[utf8]{inputenc}
\usepackage[english]{babel}
\usepackage{geometry}
\geometry{margin=2.5cm}
\usepackage{setspace}
\usepackage{microtype}
\usepackage{siunitx}
\begin{document}
% ----------------------------------------------------------------

\begin{center}
  {\large\bfseries
    Learning Motor Control Strategies in Musculoskeletal Systems\\[4pt]
    Using Deep Reinforcement Learning:\\[2pt]
    Replacing Physics-Based Muscle Solvers with Trained Policies}\\[10pt]
  {\normalsize Mohamed Naceur Hammami 
  %\quad|\quad
   %Supervisor: Prof.\ Taysir Rezgui
   }\\[4pt]
  %{\small\itshape
   %Master's in Complex and Intelligent Systems ---
   %\'{E}cole Polytechnique de Tunisie, University of Carthage, 2026}
\end{center}

\bigskip
\noindent\rule{\textwidth}{0.6pt}
\smallskip
\begin{center}{\large\bfseries Abstract}\end{center}
\smallskip
\noindent\rule{\textwidth}{0.6pt}
\bigskip

% ---- PROBLEM STATEMENT ----------------------------------------
\noindent
Real-time musculoskeletal control requires solving the Hill fiber--tendon
equilibrium at each simulation timestep to obtain per-muscle forces. The
standard Newton--Raphson approach iterates 5--40 times per muscle, imposing
runtime computation that scales with both the number of muscles and the
control frequency --- a meaningful burden for resource-constrained embedded
rehabilitation systems. This work proposes replacing the online iterative
Newton--Raphson solver with a trained deep reinforcement learning (DRL)
policy that shifts computation entirely to an offline training phase,
enabling control via a single neural network forward pass at deployment;
the policy outputs steady-state muscle activations $a \in [0,1]^9$, which
are then passed directly to the MuJoCo physics engine.

% ---- METHOD ---------------------------------------------------
\medskip\noindent
A nine-muscle upper-limb model was implemented in MuJoCo~3.2 with a
Hill-type muscle representation including tendon compliance, pennation
dynamics, and first-order activation dynamics, with anatomical parameters
drawn from Holzbaur et~al.\ (2005) and Winter (2009). Four DRL algorithms
--- DDPG, TD3, SAC, and PPO --- were trained over $2 \times 10^{6}$ steps
with domain randomization ($\pm15$\,\% segment mass, $\pm20$\,\% actuator
gain, $\pm30$\,\% joint damping) across ten independent random seeds with
Optuna-optimized hyperparameters. Statistical significance and effect sizes
were assessed using non-parametric testing and bootstrap confidence intervals.

% ---- COMPUTATIONAL PERFORMANCE --------------------------------
\medskip\noindent
\textbf{Computational performance.}\;
Under identical framework conditions (Python/PyTorch on a host CPU), the
trained SAC policy executes in $295\,\mu\text{s}$ per control step versus
$2165\,\mu\text{s}$ for the Newton--Raphson equilibrium solver. This figure is
reported for completeness but does \emph{not} support a deployment claim: it
times forward simulation rather than load sharing, and against a direct solve
of the load-sharing problem the network is roughly an order of magnitude
slower. Critically, the policy's execution time
is \emph{deterministic and branch-free} --- it consists of two fixed matrix
multiplications regardless of the biomechanical state --- whereas the N-R
solver's loop count varies with convergence conditions, precluding hard
real-time guarantees on embedded hardware. This property makes the policy
directly deployable on neural processing units (NPUs). The parameter count
($393{,}750$) and memory footprint (\SI{1.5}{\mega\byte} single precision,
\SI{385}{\kilo\byte} at eight-bit quantisation) are likewise fixed and
independent of the biomechanical state.

% ---- TASK PERFORMANCE -----------------------------------------
\medskip\noindent
\textbf{Task performance.}\;
On a three-dimensional reaching task, the best configuration attained a mean
final error of \SI{0.106}{\meter} and a closest approach of
\SI{0.074}{\meter}, ending within \SI{5}{\centi\meter} in $34$\,\% of episodes
and within \SI{10}{\centi\meter} in $58$\,\%. Under a strict success criterion
--- within \SI{2}{\centi\meter} and stationary --- performance was
$0$--$4$\,\%, a criterion that a privileged controller with full state access
satisfies only $\approx 7$\,\% of the time on this model. No statistical
significance testing is reported: the comparison uses three seeds and the
principal analysis a single seed.

% ---- MUSCLE FORCE FIDELITY ------------------------------------
\medskip\noindent
\textbf{Load-sharing fidelity.}\;
The learned per-muscle forces were compared against classical static
optimisation, posed on the identical trajectory with the identical required
joint torques and verified to machine precision. The comparison is conditional:
static optimisation is given the policy's own movement and asked only how to
distribute force across it, so trajectory realism is not assessed. The policy reproduces the
coordination pattern of the classical solution ($r = 0.81 \pm 0.04$ over five
independent training runs; pattern cosine $0.813$ against a shuffled null of
$0.372$) but expends $1.91 \pm 0.15$ times the minimum effort needed to generate
those torques --- about 44\,\% above the $1.33$ a conventional controller
attains on the same task. That discrepancy proved to be an artefact of the
reward's weighting: the effort term carried roughly one twenty-eighth the weight
of the precision bonus, and raising it brings the ratio to $1.34$, at the
achievable floor, with agreement rising to $r = 0.932$. The excess is anatomically localised: the
three deltoids agree to within $2.5$--$10$\,\% normalised RMSE, whereas the
lateral and medial heads of triceps are driven at three to four and a half
times the prescribed force, indicating sustained co-contraction against the
biceps group. Across all fifteen trained policies the discrepancy ranges from
$1.21$ to $2.30\times$, and the ordering by muscular economy is approximately
the inverse of the ordering by reaching accuracy.

% ---- CONCLUSION -----------------------------------------------
\medskip\noindent
These results support replacing the Newton--Raphson solver with a trained
policy where the quantity of interest is which muscles act and in what
relative proportion, and do not support it where the absolute magnitude of
individual muscle forces matters. Amortising solver computation into a
one-time training phase is computationally viable; the learned solution
recovers the structure of the classical load-sharing criterion without ever
being shown it, but not its efficiency.

\bigskip
\noindent\textbf{Keywords:}
musculoskeletal model,
Hill muscle model,
Newton--Raphson solver,
deep reinforcement learning,
Soft Actor-Critic,
motor control,
muscle synergies,
domain randomization,
amortized computation,
embedded rehabilitation,
MuJoCo.

% --- End standalone document (uncomment if compiling independently) ---
\end{document}
% ----------------------------------------------------------------------
 in main doc,
%  or uncomment the preamble below for standalone compilation.
% ===========================================================

% --- Standalone preamble (uncomment to compile independently) ---
\documentclass[12pt, a4paper]{article}
\usepackage[T1]{fontenc}
\usepackage[utf8]{inputenc}
\usepackage[english]{babel}
\usepackage{geometry}
\geometry{margin=2.5cm}
\usepackage{setspace}
\usepackage{microtype}
\usepackage{siunitx}
\begin{document}
% ----------------------------------------------------------------

\begin{center}
  {\large\bfseries
    Learning Motor Control Strategies in Musculoskeletal Systems\\[4pt]
    Using Deep Reinforcement Learning:\\[2pt]
    Replacing Physics-Based Muscle Solvers with Trained Policies}\\[10pt]
  {\normalsize Mohamed Naceur Hammami 
  %\quad|\quad
   %Supervisor: Prof.\ Taysir Rezgui
   }\\[4pt]
  %{\small\itshape
   %Master's in Complex and Intelligent Systems ---
   %\'{E}cole Polytechnique de Tunisie, University of Carthage, 2026}
\end{center}

\bigskip
\noindent\rule{\textwidth}{0.6pt}
\smallskip
\begin{center}{\large\bfseries Abstract}\end{center}
\smallskip
\noindent\rule{\textwidth}{0.6pt}
\bigskip

% ---- PROBLEM STATEMENT ----------------------------------------
\noindent
Real-time musculoskeletal control requires solving the Hill fiber--tendon
equilibrium at each simulation timestep to obtain per-muscle forces. The
standard Newton--Raphson approach iterates 5--40 times per muscle, imposing
runtime computation that scales with both the number of muscles and the
control frequency --- a meaningful burden for resource-constrained embedded
rehabilitation systems. This work proposes replacing the online iterative
Newton--Raphson solver with a trained deep reinforcement learning (DRL)
policy that shifts computation entirely to an offline training phase,
enabling control via a single neural network forward pass at deployment;
the policy outputs steady-state muscle activations $a \in [0,1]^9$, which
are then passed directly to the MuJoCo physics engine.

% ---- METHOD ---------------------------------------------------
\medskip\noindent
A nine-muscle upper-limb model was implemented in MuJoCo~3.2 with a
Hill-type muscle representation including tendon compliance, pennation
dynamics, and first-order activation dynamics, with anatomical parameters
drawn from Holzbaur et~al.\ (2005) and Winter (2009). Four DRL algorithms
--- DDPG, TD3, SAC, and PPO --- were trained over $2 \times 10^{6}$ steps
with domain randomization ($\pm15$\,\% segment mass, $\pm20$\,\% actuator
gain, $\pm30$\,\% joint damping) across ten independent random seeds with
Optuna-optimized hyperparameters. Statistical significance and effect sizes
were assessed using non-parametric testing and bootstrap confidence intervals.

% ---- COMPUTATIONAL PERFORMANCE --------------------------------
\medskip\noindent
\textbf{Computational performance.}\;
Under identical framework conditions (Python/PyTorch on a host CPU), the
trained SAC policy executes in $295\,\mu\text{s}$ per control step versus
$2165\,\mu\text{s}$ for the Newton--Raphson equilibrium solver. This figure is
reported for completeness but does \emph{not} support a deployment claim: it
times forward simulation rather than load sharing, and against a direct solve
of the load-sharing problem the network is roughly an order of magnitude
slower. Critically, the policy's execution time
is \emph{deterministic and branch-free} --- it consists of two fixed matrix
multiplications regardless of the biomechanical state --- whereas the N-R
solver's loop count varies with convergence conditions, precluding hard
real-time guarantees on embedded hardware. This property makes the policy
directly deployable on neural processing units (NPUs). The parameter count
($393{,}750$) and memory footprint (\SI{1.5}{\mega\byte} single precision,
\SI{385}{\kilo\byte} at eight-bit quantisation) are likewise fixed and
independent of the biomechanical state.

% ---- TASK PERFORMANCE -----------------------------------------
\medskip\noindent
\textbf{Task performance.}\;
On a three-dimensional reaching task, the best configuration attained a mean
final error of \SI{0.106}{\meter} and a closest approach of
\SI{0.074}{\meter}, ending within \SI{5}{\centi\meter} in $34$\,\% of episodes
and within \SI{10}{\centi\meter} in $58$\,\%. Under a strict success criterion
--- within \SI{2}{\centi\meter} and stationary --- performance was
$0$--$4$\,\%, a criterion that a privileged controller with full state access
satisfies only $\approx 7$\,\% of the time on this model. No statistical
significance testing is reported: the comparison uses three seeds and the
principal analysis a single seed.

% ---- MUSCLE FORCE FIDELITY ------------------------------------
\medskip\noindent
\textbf{Load-sharing fidelity.}\;
The learned per-muscle forces were compared against classical static
optimisation, posed on the identical trajectory with the identical required
joint torques and verified to machine precision. The comparison is conditional:
static optimisation is given the policy's own movement and asked only how to
distribute force across it, so trajectory realism is not assessed. The policy reproduces the
coordination pattern of the classical solution ($r = 0.81 \pm 0.04$ over five
independent training runs; pattern cosine $0.813$ against a shuffled null of
$0.372$) but expends $1.91 \pm 0.15$ times the minimum effort needed to generate
those torques --- about 44\,\% above the $1.33$ a conventional controller
attains on the same task. That discrepancy proved to be an artefact of the
reward's weighting: the effort term carried roughly one twenty-eighth the weight
of the precision bonus, and raising it brings the ratio to $1.34$, at the
achievable floor, with agreement rising to $r = 0.932$. The excess is anatomically localised: the
three deltoids agree to within $2.5$--$10$\,\% normalised RMSE, whereas the
lateral and medial heads of triceps are driven at three to four and a half
times the prescribed force, indicating sustained co-contraction against the
biceps group. Across all fifteen trained policies the discrepancy ranges from
$1.21$ to $2.30\times$, and the ordering by muscular economy is approximately
the inverse of the ordering by reaching accuracy.

% ---- CONCLUSION -----------------------------------------------
\medskip\noindent
These results support replacing the Newton--Raphson solver with a trained
policy where the quantity of interest is which muscles act and in what
relative proportion, and do not support it where the absolute magnitude of
individual muscle forces matters. Amortising solver computation into a
one-time training phase is computationally viable; the learned solution
recovers the structure of the classical load-sharing criterion without ever
being shown it, but not its efficiency.

\bigskip
\noindent\textbf{Keywords:}
musculoskeletal model,
Hill muscle model,
Newton--Raphson solver,
deep reinforcement learning,
Soft Actor-Critic,
motor control,
muscle synergies,
domain randomization,
amortized computation,
embedded rehabilitation,
MuJoCo.

% --- End standalone document (uncomment if compiling independently) ---
\end{document}
% ----------------------------------------------------------------------
